\documentclass[11pt,a4paper]{article}
\usepackage[italian]{babel}
\usepackage[margin=2.5cm]{geometry}
\usepackage{parskip}
\usepackage{amsmath, amssymb, amsthm}
\usepackage{xcolor}
\usepackage{listings}
\usepackage{tikz}
\usepackage[most]{tcolorbox}
\usepackage{hyperref}

% --- Riquadri con bordo nero, senza numero ---
\tcbset{nota/.style={colback=white, colframe=black, boxrule=0.5pt,
  sharp corners}}
\NewTColorBox{definizione}{}{nota,
  before upper={\textbf{Definizione.}\ }}
\NewTColorBox{osservazione}{}{nota,
  before upper={\textbf{Osservazione.}\ }}
\theoremstyle{definition}
\newtheorem*{esempio}{Esempio}
\newtheorem*{esercizio}{Esercizio}

% --- Codice C ---
% Uso: \begin{codice} ... \end{codice}; \begin{risultato} ... \end{risultato}
%      \lstinline|printf("ciao");| per il codice dentro al testo
\lstdefinestyle{c}{
  language=C,
  basicstyle=\ttfamily\small,
  keywordstyle=\color{blue}\bfseries,
  commentstyle=\color{gray}\itshape,
  stringstyle=\color{red!70!black},
  numbers=left,
  numberstyle=\tiny\color{gray},
  numbersep=8pt,
  columns=flexible,
  keepspaces=true,
  breaklines=true,
  showstringspaces=false,
  tabsize=4,
  morekeywords={include, define},
  literate={à}{{\`a}}1 {è}{{\`e}}1 {é}{{\'e}}1 {ì}{{\`i}}1
           {ò}{{\`o}}1 {ù}{{\`u}}1
}
\lstset{style=c}
\newtcblisting{codice}{listing only, nota, left=6mm,
  listing options={style=c}}
\newtcblisting{risultato}{listing only, colback=black!5, colframe=black,
  boxrule=0.5pt, sharp corners,
  listing options={basicstyle=\ttfamily\small, numbers=none, language={}}}

% --- Diagrammi di flusso (simboli standard) ---
\usetikzlibrary{shapes.geometric, arrows.meta, positioning}
\tikzset{
  terminale/.style = {ellipse, draw, minimum width=2.5cm, minimum height=0.9cm},
  io/.style        = {trapezium, trapezium left angle=70, trapezium right angle=110,
                      draw, minimum width=2.5cm, minimum height=0.9cm},
  processo/.style  = {rectangle, draw, minimum width=2.5cm, minimum height=0.9cm},
  decisione/.style = {diamond, draw, aspect=2, minimum width=2.5cm},
  freccia/.style   = {thick, -{Stealth}}
}

\title{Struttura di un programma C e compilazione}
\author{Marco Cerbella}
\date{Lezione 4 -- 29 settembre 2026}

\begin{document}
\maketitle

\section{Struttura di un programma}

Gli elementi costitutivi di un programma C sono le \emph{funzioni}, che possono
chiamarsi a vicenda.
\begin{itemize}
    \item Ogni funzione ha uno scopo preciso; le funzioni \emph{non} si possono
          annidare.
    \item Una funzione contiene istruzioni eseguite in sequenza; le istruzioni
          si raggruppano in \emph{blocchi}.
    \item Si possono usare funzioni della libreria standard (\lstinline|printf|)
          o scriverne di proprie.
    \item Ogni programma deve definire almeno la funzione \lstinline|main()|:
          è la prima funzione eseguita ed è il livello principale di controllo.
\end{itemize}

\section{Hello World e compilazione con gcc}

\begin{codice}
#include <stdio.h>

int main() {
    printf("Hello World\n");
    return 0;
}
\end{codice}

Salvato il codice in \texttt{test.c}, da terminale:
\begin{codice}
gcc -o test test.c     // compila e crea l'eseguibile "test"
./test                 // esegue il programma
\end{codice}

\begin{osservazione}
Per il corso e per l'esame si compila \emph{solo con gcc} da terminale (anche
se si usa un editor che ha un terminale integrato). Senza \lstinline|-o| il
file eseguibile si chiama \texttt{a.out}.
\end{osservazione}

\section{Input e output: printf e scanf}

\begin{esempio}[Somma di due interi]
\begin{codice}
#include <stdio.h>
int main() {
    int firstNumber, secondNumber, sumOfTwoNumbers;

    printf("Enter two integers: ");

    // gli interi inseriti vengono memorizzati con scanf()
    scanf("%d %d", &firstNumber, &secondNumber);

    // la somma viene salvata in sumOfTwoNumbers
    sumOfTwoNumbers = firstNumber + secondNumber;

    // stampa della somma
    printf("%d + %d = %d", firstNumber, secondNumber, sumOfTwoNumbers);

    return 0;
}
\end{codice}
\begin{risultato}
Enter two integers: 12
11
12 + 11 = 23
\end{risultato}
\end{esempio}

\begin{esempio}[Quoziente e resto]
\begin{codice}
#include <stdio.h>
int main() {
    int dividend, divisor, quotient, remainder;

    printf("Enter dividend: ");
    scanf("%d", &dividend);

    printf("Enter divisor: ");
    scanf("%d", &divisor);

    quotient = dividend / divisor;    // quoziente
    remainder = dividend % divisor;   // resto

    printf("Quotient = %d\n", quotient);
    printf("Remainder = %d", remainder);

    return 0;
}
\end{codice}
\begin{risultato}
Enter dividend: 25
Enter divisor: 4
Quotient = 6
Remainder = 1
\end{risultato}
\end{esempio}

\begin{osservazione}
\lstinline|%d| è il segnaposto per un intero; in \lstinline|scanf| le variabili
vanno precedute da \lstinline|&|.
\end{osservazione}

\section{Programmazione procedurale e modulare}

\begin{definizione}
La \emph{programmazione procedurale} è un paradigma, derivato dalla
programmazione strutturata, basato sul concetto di chiamata di procedura. Le
procedure (o routine, sottoprogrammi, funzioni) contengono una serie di passi di
calcolo da eseguire. Il C è imperativo e procedurale: le istruzioni sono
raggruppate in funzioni.
\end{definizione}

\begin{esempio}
\texttt{mainfile.c}
\begin{codice}
#include <stdio.h>
void myPrintHello();

int main() {
    myPrintHello();
    return(0);
}

void myPrintHello(void) {
    printf("Hello!\n");
    return;
}
\end{codice}
\end{esempio}

\begin{definizione}
La \emph{programmazione modulare} è una tecnica di progettazione che separa le
funzionalità di un programma in \emph{moduli} indipendenti e intercambiabili,
ciascuno contenente tutto il necessario per realizzare un solo aspetto della
funzionalità.
\end{definizione}

Ad esempio il codice si può dividere in file diversi.

\texttt{mainfile.c}:
\begin{codice}
int main() {
    myPrintHello();
    return(0);
}
\end{codice}

\texttt{hello.c}:
\begin{codice}
#include <stdio.h>

void myPrintHello(void) {
    printf("Hello!\n");
    return;
}
\end{codice}

e si compilano insieme:
\begin{codice}
gcc -o executable mainfile.c hello.c
\end{codice}
\end{document}
